% !TeX root = ..\ejemplo.tex


%%%%%%%% PREÁMBULO %%%%%%%%%%%%

\title{Plantilla para prácticas de UPIITA}
\usepackage{preambulo-plantilla}

%%%%%%%% TERMINA PREÁMBULO %%%%%%%%%%%%

\begin{document}

%%%%%%%%%%%%%%%%%%%%%%%%%%%%%%%%%% PORTADA %%%%%%%%%%%%%%%%%%%%%%%%%%%%%%%%%%%%%%%%%%%%
%%%
\begin{center}
	\newcommand{\HRule}{\rule{\linewidth}{0.5mm}}

	\begin{minipage}{0.48\textwidth} \begin{flushleft}
			\includegraphics[scale = 0.63]{Imagenes/logo_upiita}
		\end{flushleft}\end{minipage}
	\begin{minipage}{0.48\textwidth} \begin{flushright}
			\includegraphics[scale = 0.35]{Imagenes/IPN}
		\end{flushright}\end{minipage}

	%%%
	\vspace*{-1.5cm}
	%%%	
	\textsc{\huge Instituto Politécnico\\ \vspace{5px} Nacional}\\[1.5cm]

	\textsc{\LARGE Unidad Profesional Interdisciplinaria en Ingeniería y
		Tecnologías Avanzadas}\\[1.5cm]

	\begin{minipage}{0.9\textwidth}
		\begin{center}
			\textsc{\LARGE Proyecto Web}
		\end{center}
	\end{minipage}\\[0.5cm]


	\vspace*{1cm}

	\HRule \\[0.4cm]
	{ \huge \bfseries \docTitulo}\\[0.4cm]

	\HRule \\[1.5cm]

	\begin{table}[H]
		\centering
		\begin{tabular}{L{0.4\linewidth} R{0.4\linewidth}}
			\emph{Autores:} & \emph{Profesor:} \\

			{%% Autores
			Keb Sebastián Fernández Guerrero

			Ricardo Rodriguez Escalera 
			}               & {                %% Profesor
					Francisco Antonio Polanco Montelongo
			}                                  \\
		\end{tabular}
	\end{table}
	\vspace{1cm}

	\flushleft{\textbf{\Large Ing. Telemática}	}\\
	\vspace{2cm}
	\begin{center}
		{\large \today}
	\end{center}
\end{center}

\newpage
%%%

%%%%%%%%%%%%%%%%%%%% TERMINA PORTADA %%%%%%%%%%%%%%%%%%%%%%%%%%%%%%%%

\tableofcontents
\newpage

%%  Compilacion segun parametros
%% Idea de: https://tex.stackexchange.com/questions/1492/passing-parameters-to-a-document

\ifdefined\docRequerimientos
  \import{10-requerimientos}{main}
\else
  \ifdefined\docEjemplo
    \import{01-ejemplo}{main}
  \fi
\fi


\end{document}